\clearpage
\fancyhead[L]{\small BESTIMMUNG DER GENAUIGKEIT DER BEOBACHTUNGEN}
\pdfbookmark[0]{Bestimmung der Genauigkeit der Beobachtungen}{genauigkeit}
\SourceLeaf{114}
\begin{center}
\rule{.8\linewidth}{.3pt}\par\smallskip
{\small Zeitschrift für Astronomie herausgegeben von B. von \textsc{Lindenau} und J. G. F. \textsc{Bohnenberger}.\par
Erster Band. Heft für März und April 1816.\par}
\smallskip\rule{.8\linewidth}{.3pt}\par\bigskip
{\large\itshape Bestimmung der Genauigkeit der Beobachtungen.\par}
\end{center}
\renewcommand{\Article}[1]{\par\medskip\Needspace{4\baselineskip}\pdfbookmark[1]{Art. #1}{genauigkeit-art-#1}\centerline{\textbf{#1.}}\nobreak\smallskip}
\Article{1}
Bei der Begründung der sogenannten Methode der kleinsten Quadrate wird angenommen, dass die Wahrscheinlichkeit eines Beobachtungsfehlers \(\Delta\) durch die Formel
\[
\frac{h}{\surd\pi}.e^{-hh\Delta\Delta}
\]
ausgedrückt wird, wo \(\pi\) den halben Kreisumfang, \(e\) die Basis der hyperbolischen Logarithmen, auch \(h\) eine Constante bedeutet, die man nach Art. 178 der \textit{Theoria Motus Corporum Coelestium} als das Maass der Genauigkeit der Beobachtungen ansehen kann. Bei Anwendung der Methode der kleinsten Quadrate auf die Ausmittelung der wahrscheinlichsten Werthe derjenigen Grössen. von welchen die Beobachtungen abhängen, braucht man den Werth der Grösse \(h\) gar nicht zu kennen; auch das \textit{Verhältniss} der Genauigkeit der Resultate zu der Genauigkeit der Beobachtungen ist von \(h\) unabhängig. Inzwischen ist immer eine Kenntniss dieser Grösse selbst interessant und lehrreich, und ich will daher zeigen, wie man durch die Beobachtungen selbst zu einer solchen Kenntniss gelangen mag.

\Article{2}
Ich lasse zuerst einige den Gegenstand erläuternde Bemerkungen vorausgehen. Der Kürze wegen bezeichne ich den Werth des Integrals
\SourceLeaf{115}
\[
\int\frac{2e^{-tt}\,\mathrm{d}t}{\surd\pi}
\]
von \(t=0\) an gerechnet, durch \(\theta t\). Einige einzelne Werthe werden von dem Gange dieser Function eine Vorstellung geben. Man hat
\[
\begin{array}{r@{\;=\;}l@{\;=\;}l}
0{,}5000000&\theta\,0{,}4769363&\theta\rho\\
0{,}6000000&\theta\,0{,}5951161&\theta\,1{,}247790\rho\\
0{,}7000000&\theta\,0{,}7328691&\theta\,1{,}536618\rho\\
0{,}8000000&\theta\,0{,}9061939&\theta\,1{,}900032\rho\\
0{,}8427008&\theta\,1&\theta\,2{,}096716\rho\\
0{,}9000000&\theta\,1{,}1630872&\theta\,2{,}438664\rho\\
0{,}9900000&\theta\,1{,}8213864&\theta\,3{,}818930\rho\\
0{,}9990000&\theta\,2{,}3276754&\theta\,4{,}880475\rho\\
0{,}9999000&\theta\,2{,}7510654&\theta\,5{,}768204\rho\\
1&\multicolumn{2}{l}{\theta\infty}
\end{array}
\]

Die Wahrscheinlichkeit, dass der Fehler einer Beobachtung zwischen den Grenzen \(-\Delta\) und \(+\Delta\) liege, oder, ohne Rücksicht auf das Zeichen, nicht grösser als \(\Delta\) sei, ist
\[
=\int\frac{he^{-hhxx}\,\mathrm{d}x}{\surd\pi}
\]
wenn man das Integral von \(x=-\Delta\) bis \(x=+\Delta\) ausdehnt, oder doppelt so gross, wie dasselbe Integral von \(x=0\) bis \(x=\Delta\) genommen, mithin
\[
=\theta h\Delta
\]

Die Wahrscheinlichkeit, dass der Fehler nicht unter \(\frac\rho h\) sei, ist also \(=\frac12\), oder der Wahrscheinlichkeit des Gegentheils gleich: wir wollen diese Grösse \(\frac\rho h\) den \textit{wahrscheinlichen Fehler} nennen, und mit \(r\) bezeichnen. Hingegen ist die Wahrscheinlichkeit, dass der Fehler über \(2{,}438664r\) hinausgehe, nur \(\frac1{10}\); die Wahrscheinlichkeit, dass der Fehler über \(3{,}818930r\) steige, nur \(\frac1{100}\) u.\,s.\,w.

\Article{3}
Wir wollen nun annehmen, dass bei \(m\) wirklich angestellten Beobachtungen die Fehler \(\alpha,\,\beta,\,\gamma,\,\delta\) u.\,s.\,w. begangen sind, und untersuchen, was sich daraus in Beziehung auf den Werth von \(h\) und \(r\) schliessen lasse. Macht man zwei
